\documentclass[a4paper,12pt]{article}

\input{../../Preambulo.tex}
\newcommand{\assunto}{POLINÔMIOS}
\newcommand{\atividade}{lista} % prova ou lista
\newcommand{\turma}{INFORMÁTICA 3.\textordmasculine\ ANO}
\newcommand{\numProva}{A}
\newcommand{\professor}{Igor Martins Silva}
\newcommand{\bimestre}{2}
\newcommand{\ano}{2026}
\newcommand{\mesTexto}{junho}
\newcommand{\mesNum}{06}
\newcommand{\dia}{26}
\newcommand{\dataTexto}{\dia\ de \mesTexto\ de \ano}
\newcommand{\dataNun}{\dia/\mesNum/\ano}

\newcommand{\titulo}{}

\ifdefstring{\atividade}{prova}
  {\renewcommand{\titulo}{Avaliação de Matemática}}
{
\ifdefstring{\atividade}{lista}
  {\renewcommand{\titulo}{Lista de exercícios de Matemática}}
{
\ifdefstring{\atividade}{as}
  {\renewcommand{\titulo}{Avaliação Somativa}}
{
\ifdefstring{\atividade}{ad}
  {\renewcommand{\titulo}{Avaliação Diagnóstica}}
  {\renewcommand{\titulo}{Atividade de Matemática}}
}}}

\newcommand{\emtipo}{%
  \ifdefstring{\tipo}{individual}{\tipo}{em \tipo}%
}

\newcommand{\subtitulobase}[3]{%
    % #1 = título da 1ª coluna
    % #2 = conteúdo da 1ª coluna
    % #3 = mostra número da prova? (vazio ou algo)
    
    \begin{center}
        CEFET-MG -- Campus Contagem \\[3pt]
        \titulo\ -- \bimestre.\textordmasculine\ bimestre de \ano \\[3pt]
        \professor
    \end{center}
    
    \begin{center}
        \vspace{15pt}
        \begin{tabular}{|C{210pt}|C{70pt}|C{210pt}|}
            \hline
            \textbf{#1} &
            \textbf{DATA} &
            \textbf{TURMA} \\
            \hline
            #2 & \dataNun & \turma \\
            \hline
        \end{tabular}
    \end{center}
    
    #3
    
    \vspace{15pt}
}

\newcommand{\subtitulolis}{%
  \subtitulobase{ASSUNTO}{\assunto}{}%
}

\newcommand{\subtitulopr}{%
  \subtitulobase{ASSUNTO}{\assunto}{\boxed{\numProva}}%
}

\newcommand{\subtituloas}{%
  \subtitulobase{DISCIPLINA}{MATEMÁTICA}{}%
}

\newcommand{\subtituload}{%
  \subtitulobase{DISCIPLINA}{MATEMÁTICA}{}%
}

\newcommand{\valorquestao}{}

\newenvironment{exercicioBanco}[1][]{%
    \ifdefstring{\atividade}{prova}{%
        \begin{exercicio}[#1]%
    }{%
        \begin{exercicio}%
    }%
}{%
    \end{exercicio}
}

\ifdefstring{\atividade}{lista}{
    \pagecolor{background-color}
}{
    
}



\hypersetup{
    pdfauthor={\professor},
    pdftitle={\titulo \ -- \assunto \ -- \turma},
}

\title{\titulo \ -- \assunto \ -- \turma}
\author{\professor}
\date{\data}

\graphicspath{
    {../../Polinomios/}
}

%************* INÍCIO DO DOCUMENTO *************
\begin{document}
    \subtitulolis
        
    \phantomsection
    \addcontentsline{toc}{section}{Exercício 1}
    \input{../P_Exercicio1_1}
    \noindent\rule{\linewidth}{1pt}  % linha horizontal fina
    
    \phantomsection
    \addcontentsline{toc}{section}{Exercício 2}
    \input{../P_Exercicio2_1}
    \noindent\rule{\linewidth}{1pt}  % linha horizontal fina
    
    \phantomsection
    \addcontentsline{toc}{section}{Exercício 3}
    \input{../P_Exercicio3_1}
    \noindent\rule{\linewidth}{1pt}  % linha horizontal fina
    
    \phantomsection
    \addcontentsline{toc}{section}{Exercício 4}
    \input{../P_Exercicio4_1}
    \noindent\rule{\linewidth}{1pt}  % linha horizontal fina
    
    \phantomsection
    \addcontentsline{toc}{section}{Exercício 5}
    \input{../P_Exercicio5_1}
    \noindent\rule{\linewidth}{1pt}  % linha horizontal fina
    
    \phantomsection
    \addcontentsline{toc}{section}{Exercício 6}
    \input{../P_Exercicio6_1}
    \noindent\rule{\linewidth}{1pt}  % linha horizontal fina
\end{document}
%*************** FINAL DO DOCUMENTO ***************
